\documentclass[10pt,a4paper]{amsart}
\usepackage[margin=19mm]{geometry}
\usepackage{amsmath,amssymb,mathtools}
\usepackage[scr=rsfs,cal=euler]{mathalfa}
\usepackage{xfrac}
\usepackage[unicode,hidelinks,bookmarks=false]{hyperref}
\newcommand{\ie}{i.e.\ }
\newcommand{\cf}{cf.\ }
\newcommand{\resp}{respectively\ }
\newcommand{\Subsection}{\subsection}
\providecommand{\nobreakdash}{\nobreak\hbox{-}}
\setlength{\emergencystretch}{2em}
\multlinegap0pt
\usepackage{bm}
\usepackage{bbm}
\usepackage{dsfont}
\usepackage{xspace}
\usepackage{cancel}
\usepackage{mathtools}
\usepackage{amsthm}
\makeatletter
\@ifundefined{corollary}{\newtheorem{corollary}{Corollary}}{}
\@ifundefined{theorem}{\newtheorem{theorem}{Theorem}}{}
\makeatother
\makeatletter
\expandafter\ifx\csname acknowledgements\endcsname\relax
\newenvironment{acknowledgements}{\bigskip\textbf{Acknowledgements.}}{}
\fi
\newcommand{\rightarrowlim}{\mathop{\rightarrow}\limits}
\makeatother
\title[Partitions with Durfee triangles of fixed size]{Partitions with Durfee triangles of fixed size}
\author{N. Guru Sharan, Armin Straub}
\date{Source: arXiv:2507.19047v1 (CC BY 4.0)}
\begin{document}
\maketitle

\noindent\emph{Source and licence.} Partitions with Durfee triangles of fixed size. Version arXiv:2507.19047v1 (CC BY 4.0), distributed under the Creative Commons Attribution 4.0 International licence, as recorded in the arXiv licence metadata for this exact version and shown on the abstract page https://arxiv.org/abs/2507.19047v1. Full rights notice: licenses/arxiv-2507.19047-CC-BY-4.0.txt.\par\smallskip

\noindent\textit{Editorial scope.} Counted unit: the Theorem whose source label is \texttt{thm:Rk:asy} in the article (source0.tex lines 941--948, source label \texttt{thm:Rk:asy}), delivered together with the article's own complete original proof of it (source0.tex lines 950--1022, one \texttt{proof} environment of 73 lines). No mathematics is written, repaired or re-derived: statement and proof are the article's own text line for line. Required context retained uncounted: eq:qpochhammer (lines 441--443); eq:Aq:denom (lines 449--451); cor:alphad:1 (lines 648--652); thm:Fk:denom (lines 721--729); eq:phik:alpha (lines 748--751); eq:cfinite:exact (lines 922--925). Every definition, display and label of those lines is retained unchanged. Omitted from this package and not redistributed: 1--440, 444--448, 452--647, 653--720, 730--747, 752--921, 926--940, 949--949, 1023--1356. Nothing inside a retained excerpt is omitted or altered: every retained body is delivered exactly as the article's lines are written, so no omission receipt of any kind accompanies this package. Citations: the retained excerpts carry no citation command at all, so no bibliography entry of the article is transcribed and no reference is redistributed. Reference adaptation: none was needed and none was made; every reference inside the delivered unit resolves inside it, so no unresolved reference or citation is produced. Printed slips kept literal: no printed word, symbol, hypothesis or number of the counted statement or of its proof was altered. Layout and class: the article's own class is \texttt{article} with packages including geometry, amsmath, amsthm, amssymb, latexsym, xcolor, hyperref, tikz, thmtools; the wrapper is the lane's pinned \texttt{amsart} editorial wrapper at 10pt on A4, which restates the theorem-like environments the retained text needs and adds the article's own macro definitions that the retained text needs (\texttt{\textbackslash rightarrowlim}), with extra wrapper packages (bm, bbm, dsfont, xspace, cancel, mathtools). The delivered numbering is the unit's own. Assets: no retained span contains a figure environment, a graphics-inclusion command, a PDF-page inclusion, a table or a TikZ picture; no image file is copied into this package, the package declares no asset dependency and no image file is redistributed with it. Licence: version arXiv:2507.19047v1 of this article is distributed under the Creative Commons Attribution 4.0 International licence, as recorded in the arXiv licence metadata for this exact version and shown on the abstract page https://arxiv.org/abs/2507.19047v1. A full rights notice is delivered at licenses/arxiv-2507.19047-CC-BY-4.0.txt. Characters: every retained excerpt is pure ASCII, so the delivered engine, XeLaTeX with its default Latin Modern text font, reports no missing character.
\begin{equation}
  (q ; q)_d = \prod_{r = 1}^d (1 - q^r) . \label{eq:qpochhammer}
\end{equation}

  \begin{equation}
    A_d (q) = \frac{\alpha_d (q)}{(q ; q)_d} \label{eq:Aq:denom}
  \end{equation}

\begin{corollary}
  \label{cor:alphad:1}Let $\alpha_d (q) = A_d (q) (q ; q)_d$ as in
  \eqref{eq:Aq:denom}. Then, for all positive integers $d$, we have $\alpha_d
  (1) = 1$.
\end{corollary}

\begin{theorem}
  \label{thm:Fk:denom}For all positive integers $k$,
  \begin{equation}
    F_k (q) = \frac{\varphi_k (q)}{(q ; q)_k} \label{eq:Fk:denom}
  \end{equation}
  where $\varphi_k (q) \in 1 + q\mathbb{Z} [q]$ is a polynomial of degree
  $k^2$. Moreover, the coefficient of $q^{k^2}$ in $\varphi_k (q)$ is $(-
  1)^{k - 1}$.
\end{theorem}

  \begin{equation}
    \varphi_k (q) = \sum_{d = 0}^k \binom{k}{d}_q q^d \alpha_d (q) \alpha_{k -
    d} (q) \in 1 + q\mathbb{Z} [q] . \label{eq:phik:alpha}
  \end{equation}

\begin{equation}
  a (n) = \sum_{j = 1}^s \sum_{r = 0}^{m_j - 1} c_{j, r} n^r \lambda_j^n,
  \label{eq:cfinite:exact}
\end{equation}

\begin{theorem}
  \label{thm:Rk:asy}Fix a positive integer $k$. As $n \rightarrowlim \infty$,
  we have
  \begin{equation}
    R_k (n) = \frac{2^k}{k!} \frac{n^{k - 1}}{(k - 1) !} + O (n^{k - 2}) .
    \label{eq:Rk:asy:first}
  \end{equation}
\end{theorem}

\begin{proof}
  By Theorem~\ref{thm:Fk:denom}, we have
  \begin{equation*}
    \mathcal{F}_k (q) = \sum_{n = 1}^{\infty} R_k (n) q^n = q^{T_k} 
     \frac{\varphi_k (q)}{(q ; q)_k}
  \end{equation*}
  with $\varphi_k (q) \in 1 + q\mathbb{Z} [q]$. We can see from
  \eqref{eq:qpochhammer} that the roots of $(q ; q)_k$ are the roots of unity
  $\zeta$ of order $m$ with $m \leq k$ and their multiplicity is $\lfloor
  k / m \rfloor$. In other words,
  \begin{equation*}
    (q ; q)_k = (- 1)^k \prod_{m = 1}^k \Phi_m^{\lfloor k / m \rfloor} (q)
  \end{equation*}
  where $\Phi_m$ is the $m^{\text{th}}$ cyclotomic polynomial (the monic
  polynomial whose roots are precisely the $m^{\text{th}}$ primitive roots of
  unity). We write $\lambda_1 = 1$ and let $\lambda_2, \lambda_3, \ldots,
  \lambda_s$ be the other roots of unity of order up to $k$. We denote with
  $m_1, m_2, \ldots$ the corresponding multiplicities. Note that $m_1 = k$ and
  that $m_j \leq k / 2$ for $j > 1$. The exact formula
  \eqref{eq:cfinite:exact} therefore implies
  \begin{equation}
    R_k (n) = \sum_{r = 0}^{k - 1} c_{1, r} n^r + \sum_{j = 2}^s \sum_{r =
    0}^{m_j - 1} c_{j, r} n^r \lambda_j^n = \sum_{r = 0}^{k - 1} c_{1, r} n^r
    + O (n^{\lfloor k / 2 \rfloor - 1}) . \label{eq:Rk:asy:1}
  \end{equation}
  It remains to determine the leading coefficient $c_{1, k - 1}$. To this end,
  we write
  \begin{equation*}
    \mathcal{F}_k (q) = q^{T_k}  \frac{\varphi_k (q)}{(q ; q)_k} =
     \frac{\tilde{N} (q)}{(1 - q)^k \tilde{D} (q)}
  \end{equation*}
  for polynomials $\tilde{N} (q), \tilde{D} (q) \in \mathbb{Z} [q]$ with
  $\tilde{N} (1) \neq 0$ and $\tilde{D} (1) \neq 0$. Note that $\tilde{N} (q)
  = q^{T_k} \varphi_k (q)$ and $\tilde{D} (q) = (q ; q)_k / (1 - q)^k$.
  Performing part of the partial fraction expansion, we then have
  \begin{equation*}
    \mathcal{F}_k (q) = \frac{\gamma_k}{(1 - q)^k} + \frac{\gamma_{k - 1}}{(1
     - q)^{k - 1}} + \cdots + \frac{\gamma_1}{1 - q} + \frac{M (q)}{\tilde{D}
     (q)}
  \end{equation*}
  where $\gamma_j$ are constants and $M (q) \in \mathbb{Z} [q]$. Using
  \begin{equation*}
    \frac{(j - 1) !}{(1 - q)^j} = \sum_{n \geq 0} (n + j - 1) (n + j -
     2) \cdots (n + 1) q^n,
  \end{equation*}
  the coefficients $c_{1, r}$ in \eqref{eq:Rk:asy:1} can be obtained from the
  coefficients $\gamma_j$. In particular, the leading coefficient can be
  computed as
  \begin{equation*}
    c_{1, k - 1} = \frac{\gamma_k}{(k - 1) !} = \frac{1}{(k - 1) !} 
     \frac{\tilde{N} (1)}{\tilde{D} (1)} .
  \end{equation*}
  Since $\tilde{N} (q) = q^{T_k} \varphi_k (q)$, we have $\tilde{N} (1) =
  \varphi_k (1)$. On the other hand, we find that $\tilde{D} (1) = k!$ because
  \begin{equation*}
    \tilde{D} (q) = \frac{(q ; q)_k}{(1 - q)^k} = \prod_{r = 1}^k \frac{1 -
     q^r}{1 - q} = \prod_{r = 1}^k (1 + q + q^2 + \cdots + q^{r - 1}) .
  \end{equation*}
  Combined, we obtain
  \begin{equation*}
    c_{1, k - 1} = \frac{1}{(k - 1) !}  \frac{\tilde{N} (1)}{\tilde{D} (1)} =
     \frac{\varphi_k (1)}{(k - 1) !k!} .
  \end{equation*}
  By \eqref{eq:Rk:asy:1}, this proves the claimed asymptotics
  \eqref{eq:Rk:asy:first} if we can show that $\varphi_k (1) = 2^k$. Indeed,
  we can conclude from $\alpha_d (1) = 1$, established in
  Corollary~\ref{cor:alphad:1}, combined with \eqref{eq:phik:alpha} that
  \begin{equation*}
    \varphi_k (1) = \sum_{d = 0}^k \binom{k}{d} \alpha_d (1) \alpha_{k - d}
     (1) = \sum_{d = 0}^k \binom{k}{d} = 2^k,
  \end{equation*}
  as desired.
\end{proof}

\end{document}
